\section*{Podsumowanie najważniejszych własności CCC}

Jeśli $\C$ jest CCC, to:

\begin{enumerate}
\item Dla każdego obiektu $A$ mamy funktor
\[
(-)\times A \colon \C \longrightarrow \C, \qquad X \longmapsto X\times A, \qquad \bigl(X \xrightarrow{f} Y\bigr) \longmapsto \bigl(X\times A \xrightarrow{f\times \id_A} Y\times A\bigr).
\]

\item Dla każdego $A$ mamy funktor
\[
(-)^{A} \colon \C \longrightarrow \C, \qquad X \longmapsto X^{A},
\]
\[
\bigl(X \xrightarrow{f} Y\bigr) \longmapsto \bigl(X^{A} \xrightarrow{f^{A}} Y^{A}\bigr), \qquad f^{A} \stackrel{df}{=} \widetilde{f\circ \varepsilon}.
\]

\item Dla każdego $A,X,Y$ mamy bijekcję
\[
\C\bigl(X\times A,\; Y\bigr) \;\cong\; \C\bigl(X,\; Y^{A}\bigr),
\]
zadaną przez transpozycję $f \mapsto \tilde{f}$.

\item Dla każdego obiektu $X,A$ mamy ewaluację
\[
\varepsilon \colon X^{A}\times A \longrightarrow X.
\]
\end{enumerate}

\begin{stwierdzenie}
Dla każdego $f\colon X\to Y$ następujący diagram jest przemienny (naturalność $\varepsilon$):
\[
\begin{tikzcd}[column sep=large, row sep=large]
X^{A}\times A \arrow[r, "\varepsilon_{X}"] \arrow[d, "f^{A}\times \id_{A}"'] & X \arrow[d, "f"] \\
Y^{A}\times A \arrow[r, "\varepsilon_{Y}"'] & Y
\end{tikzcd}
\]
\end{stwierdzenie}

\begin{proof}
Weźmy
\[
g \colon X^{A}\times A \xrightarrow{\varepsilon} X \xrightarrow{f} Y.
\]
Z definicji $f^A = \widetilde{f\circ\varepsilon}$ wynika, że $\varepsilon\circ(f^A\times\id_A) = f\circ\varepsilon$, co jest dokładnie przemiennością kwadratu powyżej.
\end{proof}

\begin{enumerate}
\setcounter{enumi}{4}
\item Mamy strzałkę $\eta\colon X \to (X\times A)^{A}$ zdefiniowaną jako
\[
\eta \stackrel{df}{=} \widetilde{\id_{X\times A}}.
\]

\item Następujący wzór jest pomocny przy składaniu strzałek wykładnikowych: dla $f\colon X\times A\to Y\times A$, $g\colon Y\times A\to Z\times A$
\[
\widetilde{g\circ f} \;=\; (\varepsilon_{Z\times A})^{A}\circ \bigl(\widetilde{g}\times\id_{A}\bigr)^{A}\circ \widetilde{f}.
\]
\end{enumerate}

\begin{uwaga}
Ćwiczenie: uzasadnić powyższą równość.
\end{uwaga}

\section*{Wykorzystanie CCC w modelowaniu obliczeń z efektami ubocznymi}

Niech w naszej kategorii istnieje specjalny obiekt $W$ odpowiadający za stan ,,świata zewnętrznego''. Funkcja, która przyjmuje argument typu $X$ i zwraca wartość typu $Y$, jednocześnie ingerując w $W$, jest typu
\[
f \colon X\times W \longrightarrow Y\times W.
\]

Z bijekcji wynika
\[
\frac{f\colon X\times W \longrightarrow Y\times W}{\widetilde{f}\colon X \longrightarrow (Y\times W)^{W}}.
\]

Jeśli zdefiniujemy
\[
\IO(Y) \;=\; (Y\times W)^{W},
\]
to nasza funkcja jest typu $\widetilde{f}\colon X \to \IO(Y)$.

Tu $\IO = (-)^W \circ ((-)\times W)$ jest funktorem na $\C$.

\subsection*{Przykłady typów funkcji w Haskellu}

\[
\texttt{putChar} \;::\; \texttt{Char} \longrightarrow \IO\,()
\]
(wyświetlenie znaku na monitorze)

\[
\texttt{getChar} \;::\; \IO\,\texttt{Char}
\]

\[
\texttt{readFile} \;::\; \texttt{FilePath} \longrightarrow \IO\,\texttt{String}
\]
\[
\texttt{writeFile} \;::\; (\texttt{FilePath},\, \texttt{String}) \longrightarrow \IO\,()
\]
